\documentclass{article}
\usepackage{ctex}
\usepackage{amsmath}
\usepackage{tikz}
\usepackage{unicode-math}
\usetikzlibrary{arrows.meta, decorations.markings,patterns}

\title{薄金属膜电场分析}
\date{\today}
\author{}
\begin{document}

\maketitle

\section{问题}

很薄的金属薄膜(厚度为几十个甚至几个原子层，其中自由电子有限)在静电场的作用下，分析金属导体内部的场强

\section{建模}
仅从经典力学进行讨论，不考虑量子力学，进行建模

\subsection{问题重述与假设}

金属薄膜很薄，厚度 $ d $ 为纳米量级（例如几个到几十个原子层，假设 $ d \approx 1 \,\rm{nm} \sim 10 \,\rm{nm} $）。

金属内部有自由电子，但数量有限（但薄膜很薄，总电子数不多）。

只考虑经典静电学（不涉及量子尺寸效应、能带结构等）

静电场由外部施加，例如上下极板加电压，薄膜置于其中。

金属在静电情况下内部电场应为零（理想导体），但自由电子有限意味着：薄膜太薄时，电荷完全屏蔽外场需要一定的面电荷密度，如果自由电子总数不足以提供所需的面电荷密度，则屏蔽不完全，内部会有电场

\subsection{建立模型}

设薄膜在 $ z $ 方向，厚度 $ d $，面积 $ A $ 很大（$ A \gg d^2 $），可视为无限大平板。

外部电场垂直于膜面（沿 $ z $ 方向），未加膜时均匀外场大小为 $ E_\text{out} $。

金属膜放入后，自由电子会移动，在上下表面形成面电荷密度 $ \sigma $ 和 $ -\sigma $，从而产生一个内部极化电场 $ E_p $ 抵消外场。


设金属膜单位体积的自由电子数 $ n $（体密度），膜体积 $ V = A d $

总自由电子数：$N = n A d.$

总自由电荷量：$Q_\text{max} = e  = e n A d$,这里 $ e $ 为电子电荷绝对值。

\subsection{静电屏蔽能力}

要完全屏蔽外场，需要表面电荷量 $ Q_\text{need} = \varepsilon_0 E_\text{out} A $。

如果 $ Q_\text{max} < Q_\text{need} $，则电子全部移到表面也不足以完全屏蔽，此时内部电场不为零。



设实际表面电荷密度为 $ \sigma_\text{real} $，由自由电子全部移到表面形成：
$$\sigma_\text{real} = \frac{Q_\text{max}}{A} = e \times n \times d.$$
这个面电荷产生的内部电场：
$$E_p = -\frac{\sigma_\text{real}}{\varepsilon_0} = -\frac{e \times n \times d}{\varepsilon_0}.$$
因此内部电场为：
$$E_\text{in} = E_\text{out} + E_p = E_\text{out} - \frac{e \times n \times d}{\varepsilon_0}.$$


\section{理论分析}

\subsection{完全屏蔽所需的面电荷密度与电荷量}

理想导体完全屏蔽外场所需的面电荷密度：
$\sigma_{\text{need}} = \varepsilon_0 E_{\text{out}}$

所需总电荷量（单面）：
$Q_{\text{need}} = \varepsilon_0 E_{\text{out}} A$
其中 $ A $ 为薄膜面积。



\subsection{薄膜可提供的最大面电荷密度}

薄膜中全部自由电子移到表面时，最大面电荷密度：
$$
\sigma_{\max} = e n d
$$
对应最大电荷量（单面）：
$$
Q_{\max} = e n A d
$$




\subsection{表面电荷密度}

$$\sigma=\begin{cases}\varepsilon_0E_{\mathrm{ext}},E_{\mathrm{ext}}\leq E_c\\\\end,E_{\mathrm{ext}}>E_c&\end{cases}$$

\subsection{内电场分布与临界电场}

如果 $ \frac{e \times n \times d} {\varepsilon_0} \ge E_\text{out} $，则 $ E_\text{in} \le 0 $（实际上电子不会全部移到一边，会调整到内部电场为零，即理想屏蔽）。
如果 $ \frac{\rm{e n d}}{ \varepsilon_0 < E_\text{out}} $，则电子全部移到表面后，内部电场仍为正（沿外场方向），大小为：
$$E_\text{in} = E_\text{out} - \frac{e n d}{\varepsilon_0} > 0.$$
此时表面电荷密度就是最大可能值 $ \rm{e n d} $，屏蔽不完全。

当外场超过临界值  
   $$E_c = \frac{e n d}{\varepsilon_0}$$
   时，自由电子全部移到表面也无法完全屏蔽，内部存在剩余电场  
   $$E_\text{in} = E_\text{out} - E_c > 0.$$

   临界场强 $ E_c $ 与厚度 $ d $ 成正比。对于原子层级别的薄膜，$ E_c $ 可能只有 $ 10^2 \sim 10^3 \,\text{V/m} $ 量级，日常较强的静电场（如 $ 10^4 \,\text{V/m} $ 以上）就可能出现屏蔽不完全。

若 $ E_\text{out} < E_c $，则行为类似理想导体，内部电场为零，表面电荷密度为 $ \varepsilon_0 E_\text{out} < e n d $，电子未完全移到表面。



$$E_\text{in} = \max\left(0, \; E_\text{out} - \frac{e n d}{\varepsilon_0} \right)$$

$$E_\text{in}=\begin{cases}0\quad E_\text{out} \leq \frac{e n d}{\varepsilon_0} \\\ E_\text{out} - \frac{e n d}{\varepsilon_0}\quad E_\text{out} > \frac{e n d}{\varepsilon_0}&\end{cases}$$










\end{document}

